%Daten zur Institution

%\input{Dozentdaten}

%\renewcommand{\fachbereich}{Fachbereich}

%\renewcommand{\dozent}{Prof. Dr. . }

%Klausurdaten

\renewcommand{\klausurgebiet}{ }

\renewcommand{\klausurtyp}{ }

\renewcommand{\klausurdatum}{. 20}

\klausurvorspann {\fachbereich} {\klausurdatum} {\dozent} {\klausurgebiet} {\klausurtyp}

%Daten für folgende Punktetabelle

\renewcommand{\aeins}{  3  }

\renewcommand{\azwei}{  3   }

\renewcommand{\adrei}{  5  }

\renewcommand{\avier}{  2  }

\renewcommand{\afuenf}{  4  }

\renewcommand{\asechs}{  8  }

\renewcommand{\asieben}{  4  }

\renewcommand{\aacht}{  5  }

\renewcommand{\aneun}{  3  }

\renewcommand{\azehn}{  2  }

\renewcommand{\aelf}{  6  }

\renewcommand{\azwoelf}{  7  }

\renewcommand{\adreizehn}{  4   }

\renewcommand{\avierzehn}{  8  }

\renewcommand{\afuenfzehn}{  64  }

\renewcommand{\asechzehn}{    }

\renewcommand{\asiebzehn}{    }

\renewcommand{\aachtzehn}{    }

\renewcommand{\aneunzehn}{    }

\renewcommand{\azwanzig}{    }

\renewcommand{\aeinundzwanzig}{    }

\renewcommand{\azweiundzwanzig}{    }

\renewcommand{\adreiundzwanzig}{    }

\renewcommand{\avierundzwanzig}{    }

\renewcommand{\afuenfundzwanzig}{    }

\renewcommand{\asechsundzwanzig}{    }

\punktetabellevierzehn

\klausurnote

\newpage

\setcounter{section}{0}

__NOEDITSECTION__

\inputaufgabeklausurloesung
{3}

{

Definiere die folgenden
\zusatzklammer {kursiv gedruckten} {} {} Begriffe.
\aufzaehlungsechs{Die
\stichwort {zweite tangentiale Ableitung} {}
\zusatzklammer {oder die
\stichwort {tangentiale Beschleunigung} {}} {} {}
einer
\zusatzklammer {als Abbildung nach $\R^n$} {} {}
zweimal
\definitionsverweis {stetig differenzierbaren Kurve}{}{}
\maabbdisp {\gamma} {I} {Y
} {}
in eine
\definitionsverweis {differenzierbare Hyperfläche}{}{}

\mavergleichskette
{\vergleichskette
{ Y
}
{ \subseteq }{ \R^n
}
{  }{
}
{    }{
}
{   }{
}
}
{}{}{.}

}{Ein \stichwort {zusammenhängender} {} topologischer Raum $X$.

}{Die \stichwort {Tangentialabbildung in einem Punkt} {}

\mavergleichskette
{\vergleichskette
{  P
}
{ \in }{ M
}
{  }{}
{    }{}
{   }{}
}
{}{}{}
zu einer differenzierbaren Abbildung
\maabbdisp {\varphi} { M } { N
} {,}
wobei
\mathkor {} {M} {und} {N} {}
differenzierbare Mannigfaltigkeiten sind.

}{Das
\stichwort {Produkt} {}
von differenzierbaren Mannigfaltigkeiten
\mathkor {} {M} {und} {N} {.}

}{Das
\stichwort {Volumenmaß} {}
zu einer positiven Volumenform $\omega$ auf einer $n$-dimensionalen differenzierbaren Mannigfaltigkeit.

}{Eine
\stichwort {geodätische Kurve} {}
auf einer riemannschen Mannigfaltigkeit $M$.
}

}
{

\aufzaehlungsechs{Man nennt die Abbildung
\maabbeledisp {} {I} { TY
} {t} { \pi_P( \gamma^{\prime \prime} (t))
} {,}
wobei $\pi_P$ die orthogonale Projektion
\maabbdisp {} {\R^n } { T_PY
} {}
bezeichnet, die
zweite tangentiale Ableitung
\zusatzklammer {oder die
\definitionswort {tangentiale Beschleunigung}{}} {} {}
von $\gamma$.
}{Ein
\definitionsverweis {topologischer Raum}{}{}
$X$ heißt zusammenhängend, wenn es in $X$ genau zwei Teilmengen gibt
\zusatzklammer {nämlich $\emptyset$ und der Gesamtraum

\mavergleichskettek
{\vergleichskettek
{  X
}
{ \neq }{ \emptyset
}
{  }{
}
{    }{
}
{   }{
}
}
{}{}{}} {} {,}
die sowohl
\definitionsverweis {offen}{}{}
als auch
\definitionsverweis {abgeschlossen}{}{}
sind.
}{Unter der Tangentialabbildung im Punkt $P$ versteht man die Abbildung
\maabbeledisp {} {T_PM} {T_{\varphi(P)}N
} {[\gamma]} {[\varphi \circ \gamma]
} {.}
}{Es seien
\mathkor {} {(U_i,U_i',\alpha_i, i \in I)} {und} {(V_j,V_j',\beta_j, j \in J)} {}
die
\definitionsverweis {Atlanten}{}{}
von
\mathkor {} {M} {und} {N} {.}
Dann nennt man den
\definitionsverweis {Produktraum}{}{}

\mathl{M \times N}{} versehen mit den
\definitionsverweis {Karten}{}{}
\maabbdisp {\alpha_i \times \beta_j} {U_i \times V_j } { U_i' \times V_j'
} {}
\zusatzklammer {mit \mathlk{(i,j) \in I \times J}{} und \mathlk{U_i' \times V_j' \subseteq \R^m \times \R^n}{}} {} {}
das Produkt der Mannigfaltigkeiten
\mathkor {} {M} {und} {N} {.}
}{Für eine
\definitionsverweis {Borel-Menge}{}{}

\mavergleichskette
{\vergleichskette
{ T
}
{ \subseteq }{ M
}
{  }{
}
{    }{
}
{   }{
}
}
{}{}{}
wird das Maß von $T$ zu $\omega$ über eine
\definitionsverweis {abzählbare}{}{}
\definitionsverweis {Zerlegung}{}{}

\mavergleichskette
{\vergleichskette
{ T
}
{ = }{  \biguplus_{i \in I} T_i
}
{  }{
}
{    }{
}
{   }{
}
}
{}{}{}
\zusatzklammer {wobei

\mavergleichskettek
{\vergleichskettek
{ T_i
}
{ \subseteq }{U_i
}
{  }{
}
{    }{
}
{   }{
}
}
{}{}{}
ein offenes Kartengebiet und

\mavergleichskettek
{\vergleichskettek
{  \alpha_{i*} \omega
}
{ = }{ f_i dx_1 \wedge \ldots \wedge dx_n
}
{  }{
}
{    }{
}
{   }{
}
}
{}{}{}
ist} {} {}

\mavergleichskettedisp
{\vergleichskette
{
\int_{ T  }  \omega
}
{ =} { \sum_{i \in I} \int_{ \alpha_i(T_i)  }  f_i   \,  d \lambda^n
}
{ } {
}
{ } {
}
{ } {
}
}
{}{}{}
definiert.
}{Man nennt eine
\definitionsverweis {differenzierbare Kurve}{}{}
\maabbdisp {\gamma} {I} {M
} {}
eine
geodätische Kurve,
wenn

\mavergleichskettedisp
{\vergleichskette
{ \nabla_{d/dt}  \gamma'
}
{ =} { 0
}
{ } {
}
{ } {
}
{ } {
}
}
{}{}{}
auf $I$ ist.
}

 }

__NOEDITSECTION__

\inputaufgabeklausurloesung
{3}

{

Formuliere die folgenden Sätze.
\aufzaehlungdrei{Der
\stichwort {Satz über die Krümmung auf einen implizit gegebenen ebenen Kurve} {}

\mavergleichskette
{\vergleichskette
{  Y
}
{ \subseteq }{  \R^2
}
{  }{
}
{    }{
}
{   }{
}
}
{}{}{.}}{Die
\stichwort {Transformationsformel für positive Volumenformen auf Mannigfaltigkeiten} {.}}{Der Satz über die Gaußkrümmung und die Schnittkrümmung auf einer differenzierbaren Fläche

\mavergleichskette
{\vergleichskette
{  Y
}
{ \subseteq }{  \R^3
}
{  }{
}
{    }{
}
{   }{
}
}
{}{}{}}

}
{

\aufzaehlungdrei{\faktsituation {Es sei

\mavergleichskette
{\vergleichskette
{ W
}
{ \subseteq }{ \R^2
}
{  }{
}
{    }{
}
{   }{
}
}
{}{}{}
\definitionsverweis {offen}{}{,}
\maabb {h} { W } { \R
} {}
eine zweifach
\definitionsverweis {stetig differenzierbare Funktion}{}{}
und

\mavergleichskette
{\vergleichskette
{ Y
}
{ = }{ h^{-1}(c)
}
{  }{
}
{    }{
}
{   }{
}
}
{}{}{}
die
\definitionsverweis {Faser}{}{}
zu

\mavergleichskette
{\vergleichskette
{ c
}
{ \in }{ \R
}
{  }{
}
{    }{
}
{   }{
}
}
{}{}{,}
wobei $h$ in jedem Punkt von $Y$
\definitionsverweis {regulär}{}{}
sei. Es sei
\maabbdisp {N} { U } {\R^2
} {}
ein auf einer offenen Umgebung

\mavergleichskette
{\vergleichskette
{ Y
}
{ \subseteq }{ U
}
{  }{
}
{    }{
}
{   }{
}
}
{}{}{}
definiertes
\definitionsverweis {Einheitsnormalenfeld}{}{}
zu $Y$ und sei

\mavergleichskette
{\vergleichskette
{ P
}
{ \in }{ Y
}
{  }{
}
{    }{
}
{   }{
}
}
{}{}{.}}
\faktfolgerung {Dann ist für jeden tangentialen Vektor

\mavergleichskette
{\vergleichskette
{ v
}
{ \in }{ T_PY
}
{  }{
}
{    }{
}
{   }{
}
}
{}{}{}

\mavergleichskettedisp
{\vergleichskette
{ \left(  D N   \right)_{ P } \left(  v  \right)
}
{ =} { - \kappa(P) v
}
{ } {
}
{ } {
}
{ } {
}
}
{}{}{,}
wobei $\kappa(P)$ die
\definitionsverweis {Krümmung}{}{}
einer Bogenparametrisierung von $Y$ ist, die mit der durch $v, N(P)$ gegebenen Orientierung übereinstimmt.}
\faktzusatz {}
\faktzusatz {}}{\faktsituation {Es sei $M$ eine
\definitionsverweis {differenzierbare Mannigfaltigkeit}{}{}
mit
\definitionsverweis {abzählbarer Basis der Topologie}{}{}
und sei  $\omega$ eine
\definitionsverweis {positive Volumenform}{}{}
auf $M$. Es sei
\maabbdisp {\varphi} { L } { M
} {}
ein
\definitionsverweis {Diffeomorphismus}{}{}
mit der Mannigfaltigkeit $L$ und

\mavergleichskette
{\vergleichskette
{ S
}
{ \subseteq }{ L
}
{  }{
}
{    }{
}
{   }{
}
}
{}{}{}
eine
\definitionsverweis {messbare Teilmenge}{}{.}}
\faktfolgerung {Dann ist

\mavergleichskettedisp
{\vergleichskette
{ \int_S \varphi^* \omega
}
{ =} { \int_{ \varphi(S)} \omega
}
{ } {
}
{ } {
}
{ } {
}
}
{}{}{.}}
\faktzusatz {}
\faktzusatz {}}{\faktsituation {Es sei

\mavergleichskette
{\vergleichskette
{ W
}
{ \subseteq }{ \R^3
}
{  }{
}
{    }{
}
{   }{
}
}
{}{}{}
\definitionsverweis {offen}{}{}
und sei
\maabb {h} { W } { \R
} {}
eine zweifach
\definitionsverweis {stetig differenzierbare Funktion}{}{}
und

\mavergleichskette
{\vergleichskette
{ Y
}
{ = }{ h^{-1}(c)
}
{  }{
}
{    }{
}
{   }{
}
}
{}{}{}
die Faser zu

\mavergleichskette
{\vergleichskette
{ c
}
{ \in }{ \R
}
{  }{
}
{    }{
}
{   }{
}
}
{}{}{,}
wobei $h$ in jedem Punkt von $Y$
\definitionsverweis {regulär}{}{}
sei.}
\faktfolgerung {Dann ist die
\definitionsverweis {Gaußkrümmung}{}{}
von $Y$ gleich der
\definitionsverweis {Schnittkrümmung}{}{}
von $Y$.}
\faktzusatz {}
\faktzusatz {}}

 }

__NOEDITSECTION__

\inputaufgabeklausurloesung
{5}

{

Es sei

\mavergleichskette
{\vergleichskette
{ G
}
{ \subseteq }{ \R^n
}
{  }{
}
{    }{
}
{   }{
}
}
{}{}{}
\definitionsverweis {offen}{}{}
und
\maabbdisp {\varphi} { G } { \R^m
} {}
eine
\definitionsverweis {stetig differenzierbare Abbildung}{}{,}
die im Punkt

\mavergleichskette
{\vergleichskette
{ P
}
{ \in }{ G
}
{  }{
}
{    }{
}
{   }{
}
}
{}{}{}
ein surjektives
\definitionsverweis {totales Differential}{}{}
besitze. Es sei

\mavergleichskette
{\vergleichskette
{ v
}
{ \in }{ T_PZ
}
{  }{
}
{    }{
}
{   }{
}
}
{}{}{}
ein Vektor des
\definitionsverweis {Tangentialraumes an die Faser}{}{}
$Z$ zu $\varphi$ durch $P$. Zeige, dass es eine
\definitionsverweis {stetig differenzierbare Kurve}{}{}
\maabbdisp {\gamma} { ]- \epsilon, \epsilon[ } { Z
} {}
\zusatzklammer {für ein geeignetes

\mavergleichskettek
{\vergleichskettek
{ \epsilon
}
{ > }{ 0
}
{  }{
}
{    }{
}
{   }{
}
}
{}{}{}} {} {}
mit

\mavergleichskette
{\vergleichskette
{ \gamma(0)
}
{ = }{ P
}
{  }{
}
{    }{
}
{   }{
}
}
{}{}{}
und mit

\mavergleichskettedisp
{\vergleichskette
{ \gamma'(0)
}
{ =} { v
}
{ } {
}
{ } {
}
{ } {
}
}
{}{}{}
gibt.

}
{

Es ist

\mavergleichskette
{\vergleichskette
{v
}
{ \in }{ T_PZ
}
{ = }{ \operatorname{kern}  \left(D\varphi\right)_{P}
}
{    }{
}
{   }{
}
}
{}{}{.}
Nach
dem [[Satz über implizite Abbildungen/R/Fakt|Satz über implizite Abbildungen]]
gibt es eine offene Menge

\mathbed {P \in W} {}
 {W \subseteq G} {}
 {} {} {} {,}
eine offene Menge

\mavergleichskette
{\vergleichskette
{ V
}
{ \subseteq }{ \R^{n-m}
}
{  }{
}
{    }{
}
{   }{
}
}
{}{}{}
und eine stetig differenzierbare Abbildung
\maabbdisp {\psi} { V } { W
} {}
derart, dass

\mavergleichskette
{\vergleichskette
{ \psi(V)
}
{ \subseteq }{ Z \cap W
}
{  }{
}
{    }{
}
{   }{
}
}
{}{}{}
ist und $\psi$ eine
\definitionsverweis {Bijektion}{}{}
\maabbdisp {\psi} { V } { Z \cap W
} {}
induziert. Es sei

\mavergleichskette
{\vergleichskette
{ Q
}
{ \in }{ V
}
{  }{
}
{    }{
}
{   }{
}
}
{}{}{}
der Punkt mit

\mavergleichskette
{\vergleichskette
{ \psi(Q)
}
{ = }{ P
}
{  }{
}
{    }{
}
{   }{
}
}
{}{}{.}
Die Abbildung $\psi$ ist in jedem Punkt

\mavergleichskette
{\vergleichskette
{ Q
}
{ \in }{ V
}
{  }{
}
{    }{
}
{   }{
}
}
{}{}{}
\definitionsverweis {regulär}{}{}
und für das
\definitionsverweis {totale Differential}{}{}
von $\psi$ gilt

\mavergleichskettedisp
{\vergleichskette
{ \left(D \varphi \right)_{ P }  \circ \left(D \psi \right)_{ Q }
}
{ =} { 0
}
{ } {
}
{ } {
}
{ } {
}
}
{}{}{,}
also

\mavergleichskettedisp
{\vergleichskette
{ \operatorname{bild}  \left(D \psi\right)_{ Q }
}
{ =} { \operatorname{kern}  \left(D\varphi \right)_{ P }
}
{ =} { T_P Z
}
{ } {
}
{ } {
}
}
{}{}{.}
Wegen der Regularität von $\psi$ in $Q$ ist
\maabbdisp {\left(D \psi \right)_{ Q }} { \R^{n-m}  } { \R^n
} {}
injektiv und
\maabbdisp {\left(D \psi \right)_{ Q }} { \R^{n-m}  } { T_PZ
} {}
bijektiv. Es sei

\mavergleichskette
{\vergleichskette
{ u
}
{ \in }{ \R^{n-m}
}
{  }{
}
{    }{
}
{   }{
}
}
{}{}{}
das Urbild von $v$ und sei
\maabbeledisp {\theta} { ]- \epsilon, \epsilon[ } { \R^{n-m}
} { t } { Q+tu
} {,}
wobei $\epsilon$ hinreichend klein gewählt sei, dass das Bild ganz in $V$ liegt. Dann besitzt
\maabbdisp {\gamma = \psi \circ \theta} { ]- \epsilon, \epsilon[  } { \R^n
} {}
die Eigenschaft

\mavergleichskettedisp
{\vergleichskette
{ \gamma(0)
}
{ =} { \psi( \theta (0))
}
{ =} { \psi(Q)
}
{ =} { P
}
{ } {
}
}
{}{}{}
und

\mavergleichskettedisp
{\vergleichskette
{ \gamma' (0)
}
{ =} { \left(D \psi\right)_{Q} ( \theta' (0))
}
{ =} { \left(D \psi\right)_{Q} ( u)
}
{ =} { v
}
{ } {
}
}
{}{}{.}

 }

__NOEDITSECTION__

\inputaufgabeklausurloesung
{2}

{

Es seien
\mathkor {} {r} {und} {R} {}
positive reelle Zahlen mit

\mavergleichskette
{\vergleichskette
{ 0
}
{ < }{ r
}
{ < }{ R
}
{    }{
}
{   }{
}
}
{}{}{.}
Bestimme ein
\definitionsverweis {Einheitsnormalenfeld}{}{}
für den
\zusatzklammer {eingebetteten} {} {}
Torus

\mavergleichskettedisp
{\vergleichskette
{ T
}
{ =} { \left\{ (x,y,z) \in \R^3  \mid   \left(  \sqrt{x^2+y^2} -R  \right)^2 +z^2  = r^2         \right\}
}
{ } {
}
{ } {
}
{ } {
}
}
{}{}{.}

}
{

Es ist

\mavergleichskettedisp
{\vergleichskette
{ h(x,y,z)
}
{ =} { x^2+y^2 -2 \sqrt{x^2+y^2}R+z^2 + R^2
}
{ } {
}
{ } {
}
{ } {
}
}
{}{}{}
und

\mavergleichskette
{\vergleichskette
{ T
}
{ = }{ h^{-1} (r^2)
}
{  }{
}
{    }{
}
{   }{
}
}
{}{}{.}
Der Gradient von $h$ ist

\mavergleichskettedisp
{\vergleichskette
{
\operatorname{Grad} \,  h
}
{ =} { 2 \begin{pmatrix}   x   -  xR \left(  x^2+y^2  \right)^{- { \frac{   1  }{  2 } } }  \\ y   -  y R \left(  x^2+y^2  \right)^{- { \frac{   1  }{  2 } } } \\ z   \end{pmatrix}
}
{ } {
}
{ } {
}
{ } {
}
}
{}{}{.}
Das
\definitionsverweis {Einheitsnormalenfeld}{}{}
ist daher

\mavergleichskettedisp
{\vergleichskette
{ N(x,y,z)
}
{ =} { { \frac{   1  }{  \sqrt{ \left(   x   -  xR \left(  x^2+y^2  \right)^{- { \frac{   1  }{  2 } } }  \right)^2 + \left(  y   -  y R \left(  x^2+y^2  \right)^{- { \frac{   1  }{  2 } } }  \right)^2   +z^2   }  } } \begin{pmatrix}   x   -  xR \left(  x^2+y^2  \right)^{- { \frac{   1  }{  2 } } }  \\ y   -  y R \left(  x^2+y^2  \right)^{- { \frac{   1  }{  2 } } } \\ z   \end{pmatrix}
}
{ } {
}
{ } {
}
{ } {
}
}
{}{}{.}

 }

__NOEDITSECTION__

\inputaufgabeklausurloesung
{4}

{

Man erläutere die Relevanz des Satzes über implizite Abbildungen für den Aufbau der Theorie der Mannigfaltigkeiten.

}
{  }

__NOEDITSECTION__

\inputaufgabeklausurloesung
{8}

{

Beweise den Satz über die Selbstadjungiertheit der Weingartenabbildung.

}
{

Für Vektoren

\mavergleichskette
{\vergleichskette
{ v,w
}
{ \in }{ T_PY
}
{  }{
}
{    }{
}
{   }{
}
}
{}{}{}
ist

\mavergleichskettedisp
{\vergleichskette
{ \left\langle  v  ,  L_P(w)  \right\rangle
}
{ =} { \left\langle L_P(v) ,  w  \right\rangle
}
{ } {
}
{ } {
}
{ } {
}
}
{}{}{}
zu zeigen. Mit

\mavergleichskette
{\vergleichskette
{ N(P)
}
{ = }{ { \frac{
\operatorname{Grad} \,  h  (  P )  }{  \Vert {
\operatorname{Grad} \,  h  (  P ) } \Vert  } }
}
{  }{
}
{    }{
}
{   }{
}
}
{}{}{}
ist gemäß
[[Totale Differenzierbarkeit/R/Produktregel/Fakt|Lemma 45.10 (Analysis (Osnabrück 2021-2023))]]

\mavergleichskettealignhandlinks
{\vergleichskettealignhandlinks
{ \left\langle  v  ,  L_P(w)  \right\rangle
}
{ =} { -\left\langle  v  ,  \left(  D_{w}  N   \right) \left(   P   \right)   \right\rangle
}
{ =} { -\left\langle  v  ,  \left(  D_{w}  { \frac{
\operatorname{Grad} \,  h  }{  \Vert {
\operatorname{Grad} \,  h } \Vert  } }   \right) \left(   P   \right)     \right\rangle
}
{ =} { -\left\langle  v  ,  \left(  D_{w}  { \frac{   1  }{  \Vert {
\operatorname{Grad} \,  h } \Vert  } }   \right) \left(   P   \right) \cdot
\operatorname{Grad} \,  h  (  P ) + { \frac{   1  }{  \Vert {
\operatorname{Grad} \,  h  (  P ) } \Vert  } } \cdot  \left(  D_{w}
\operatorname{Grad} \,  h   \right) \left(   P   \right)   \right\rangle
}
{ =} { - { \frac{   1  }{  \Vert {
\operatorname{Grad} \,  h  (  P ) } \Vert  } } \left\langle  v  ,  \left(  D_{w}
\operatorname{Grad} \,  h   \right) \left(   P   \right)    \right\rangle
}
}
{}
{}{,}
da der erste Summand senkrecht auf dem Tangentialvektor $v$ steht. Mit Koordinatenfunktionen ist

\mavergleichskettealignhandlinks
{\vergleichskettealignhandlinks
{ \left(  D_{w}
\operatorname{Grad} \,  h   \right) \left(   P   \right)
}
{ =} { \left(  D_{w}  \left(  { \frac{   \partial   h   }{  \partial   x_1   } }   , \,   \ldots  , \,   { \frac{   \partial   h   }{  \partial   x_n   } }                 \right)    \right) \left(   P   \right)
}
{ =} { \sum_{j = 1}^n  w_j  { \frac{   \partial    }{  \partial   x_j   } }  \left(  { \frac{   \partial   h   }{  \partial   x_1   } }   , \,   \ldots  , \,   { \frac{   \partial   h   }{  \partial   x_n   } }                 \right) (P)
}
{ =} { \left(  \sum_{j = 1}^n  w_j  { \frac{   \partial    }{  \partial   x_j   } }  { \frac{   \partial   h   }{  \partial   x_1   } } (P)   , \,   \ldots  , \,   \sum_{j = 1}^n  w_j  { \frac{   \partial    }{  \partial   x_j   } }  { \frac{   \partial   h   }{  \partial   x_n   } } (P)                 \right)
}
{ } {
}
}
{}
{}{.}
Der obige Ausdruck ist somit gleich

\mavergleichskettealign
{\vergleichskettealign
{ \left\langle  v  ,  L_P(w)  \right\rangle
}
{ =} { - { \frac{   1  }{  \Vert {
\operatorname{Grad} \,  h  (  P ) } \Vert  } } \left\langle  v  ,  \left(  D_{w}
\operatorname{Grad} \,  h   \right) \left(   P   \right)   \right\rangle
}
{ =} { -  { \frac{   1  }{  \Vert {
\operatorname{Grad} \,  h  (  P ) } \Vert  } }  \sum_{i,j = 1}^n v_i  w_j  { \frac{   \partial    }{  \partial   x_j   } } { \frac{   \partial   h   }{  \partial   x_i   } }  (P)
}
{ } {
}
{ } {
}
}
{}
{}{.}
Nach
[[Differenzierbarkeit/Satz von Schwarz/Fakt|Satz 44.10 (Analysis (Osnabrück 2021-2023))]]
kann man die Reihenfolge der partiellen Ableitungen vertauschen, sodass man auch die Rollen von
\mathkor {} {v} {und} {w} {}
vertauschen kann.

 }

__NOEDITSECTION__

\inputaufgabeklausurloesung
{4}

{

Es sei $X$ ein
\definitionsverweis {kompakter Raum}{}{}
und es sei

\mavergleichskette
{\vergleichskette
{Y
}
{ \subseteq }{ X
}
{  }{
}
{    }{
}
{   }{
}
}
{}{}{}
eine
\definitionsverweis {abgeschlossene Teilmenge}{}{,}
die die
\definitionsverweis {induzierte Topologie}{}{}
trage. Zeige, dass $Y$ ebenfalls kompakt ist.

}
{

Es sei

\mavergleichskettedisp
{\vergleichskette
{ Y
}
{ =} { \bigcup_{i \in I} V_i
}
{ } {
}
{ } {
}
{ } {
}
}
{}{}{}
eine Überdeckung mit offenen Teilmenge

\mavergleichskette
{\vergleichskette
{ V_i
}
{ \subseteq }{ Y
}
{  }{
}
{    }{
}
{   }{
}
}
{}{}{.}
Dies bedeutet, dass es offene Mengen

\mavergleichskette
{\vergleichskette
{ U_i
}
{ \subseteq }{ X
}
{  }{
}
{    }{
}
{   }{
}
}
{}{}{}
gibt mit

\mavergleichskette
{\vergleichskette
{ V_i
}
{ = }{ Y \cap U_i
}
{  }{
}
{    }{
}
{   }{
}
}
{}{}{.}
Daher ist

\mavergleichskettedisp
{\vergleichskette
{ Y
}
{ \subseteq} { \bigcup_{i \in I} U_i
}
{ } {
}
{ } {
}
{ } {
}
}
{}{}{.}
Wegen der Abgeschlossenheit von $Y$ in $X$ ist $X \setminus Y$ offen und daher ist

\mavergleichskettedisp
{\vergleichskette
{ X
}
{ =} { Y \cup (X \setminus Y)
}
{ =} { \bigcup_{i \in I} U_i  \cup  (X \setminus Y)
}
{ } {
}
{ } {
}
}
{}{}{}
eine offene Überdeckung von $X$. Wegen der Kompaktheit von $X$ gibt es eine endliche Teilüberdeckung

\mavergleichskettedisp
{\vergleichskette
{ X
}
{ =} { \bigcup_{i \in J} U_i  \cup  (X \setminus Y)
}
{ } {
}
{ } {
}
{ } {}
}
{}{}{.}
Dies bedeutet wiederum

\mavergleichskettedisp
{\vergleichskette
{ Y
}
{ =} { \bigcup_{i \in J} U_i
}
{ } {
}
{ } {
}
{ } {}
}
{}{}{.}

 }

__NOEDITSECTION__

\inputaufgabeklausurloesung
{5}

{

Man gebe ein Beispiel einer
\definitionsverweis {zweidimensionalen}{}{}
\definitionsverweis {zusammenhängenden}{}{} \definitionsverweis {differenzierbaren Mannigfaltigkeit}{}{}
$M$ und einem Punkt

\mavergleichskette
{\vergleichskette
{ P
}
{ \in }{ M
}
{  }{
}
{    }{
}
{   }{
}
}
{}{}{}
derart, dass
\mathkor {} {M} {und} {M \setminus \{P\}} {}
zueinander
\definitionsverweis {diffeomorph}{}{}
sind.

}
{

Es sei

\mavergleichskettedisp
{\vergleichskette
{M
}
{ =} { \R^2 \setminus \N_+
}
{ =} { \R^2 \setminus \{ (1,0) , (2,0), (3,0), \ldots  \}
}
{ } {
}
{ } {
}
}
{}{}{,}
es werden also aus dem $\R^2$ die auf der $x$-Achse platzierten positiven natürlichen Zahlen herausgenommen. Dies ist eine offene Teilmenge im $\R^2$ und damit eine differenzierbare Mannigfaltigkeit. Diese ist wegzusammenhängend, da man beispielsweise jeden Punkt aus $M$ mit
\mathl{y \geq 0}{} durch einen geraden Weg mit
\mathl{(0,1)}{} und jeden Punkt aus $M$ mit
\mathl{y \leq 0}{} durch einen geraden Weg mit
\mathl{(0,-1)}{} verbinden kann und diese beiden Punkte ebenfalls gerade verbindbar sind. Es sei nun

\mavergleichskettedisp
{\vergleichskette
{P
}
{ =} { (0,0)
}
{ } {
}
{ } {
}
{ } {
}
}
{}{}{}
und

\mavergleichskettedisp
{\vergleichskette
{M'
}
{ =} { M \setminus \{P\}
}
{ } {
}
{ } {
}
{ } {
}
}
{}{}{.}
Die Abbildung
\maabbeledisp {} {\R^2} {\R^2
} {(x,y)} {(x-1,y)
} {,}
ist
\zusatzklammer {als Verschiebung} {} {}
ein Diffeomorphismus und das Bild von $M$ ist genau $M'$. Daher sind
\mathkor {} {M} {und} {M'=M \setminus \{P\}} {}
zueinander diffeomorph.

 }

__NOEDITSECTION__

\inputaufgabeklausurloesung
{3}

{

Man gebe für jeden
\definitionsverweis {Tangentialvektor}{}{}

\mavergleichskette
{\vergleichskette
{ v
}
{ \in }{ T_PS^2
}
{ \subseteq }{ \R^3
}
{    }{
}
{   }{
}
}
{}{}{}
mit

\mavergleichskette
{\vergleichskette
{ \Vert {v} \Vert
}
{ = }{ 1
}
{  }{
}
{    }{
}
{   }{
}
}
{}{}{}
in einem Punkt

\mavergleichskette
{\vergleichskette
{ P
}
{ \in }{ S^2
}
{ \subset }{ \R^3
}
{    }{
}
{   }{
}
}
{}{}{}
auf der
\definitionsverweis {Einheitssphäre}{}{}
einen
\definitionsverweis {differenzierbaren}{}{}
Repräsentanten
\maabbdisp {\gamma} { I } { S^2
} {}
mit

\mavergleichskette
{\vergleichskette
{ \gamma(0)
}
{ = }{ P
}
{  }{
}
{    }{
}
{   }{
}
}
{}{}{}
an.

}
{

Es sei

\mavergleichskette
{\vergleichskette
{P
}
{ = }{ \begin{pmatrix}  x  \\ y \\ z   \end{pmatrix}
}
{  }{
}
{    }{
}
{   }{
}
}
{}{}{}
und

\mavergleichskette
{\vergleichskette
{v
}
{ = }{ \begin{pmatrix}  a  \\ b \\ c   \end{pmatrix}
}
{  }{
}
{    }{
}
{   }{
}
}
{}{}{.}
Beide Vektoren sind normiert und stehen senkrecht aufeinander, da die Tangentialebene an die Sphäre senkrecht auf dem Ortsvektor steht. Wir betrachten den Weg
\maabbdisp {\gamma} {\R} { \R^3
} {}
mit

\mavergleichskettedisp
{\vergleichskette
{ \gamma(t)
}
{ =} {  \left(  \cos t  \right)  \begin{pmatrix}  x  \\ y \\ z   \end{pmatrix}  + \left(  \sin t  \right)  \begin{pmatrix}  a  \\ b \\ c   \end{pmatrix}
}
{ } {
}
{ } {
}
{ } {
}
}
{}{}{.}
Es ist

\mavergleichskettedisp
{\vergleichskette
{ \Vert { \gamma(t)} \Vert
}
{ =} { \cos^{ 2  }  t + \sin^{ 2  }  t
}
{ =} { 1
}
{ } {
}
{ } {
}
}
{}{}{,}
der Weg verläuft also ganz auf der Sphäre. Ferner ist

\mavergleichskettedisp
{\vergleichskette
{ \gamma(0)
}
{ =} { \left(  \cos  0  \right) \begin{pmatrix}  x  \\ y \\ z   \end{pmatrix} + \left(  \sin  0  \right)  \begin{pmatrix}  a  \\ b \\ c   \end{pmatrix}
}
{ =} { \begin{pmatrix}  x  \\ y \\ z   \end{pmatrix}
}
{ } {
}
{ } {
}
}
{}{}{}
und

\mavergleichskettedisp
{\vergleichskette
{ \gamma'(0)
}
{ =} { - \left(  \sin  0  \right) \begin{pmatrix}  x  \\ y \\ z   \end{pmatrix} + \left(  \cos  0  \right) \begin{pmatrix}  a  \\ b \\ c   \end{pmatrix}
}
{ =} { \begin{pmatrix}  a  \\ b \\ c   \end{pmatrix}
}
{ } {
}
{ } {
}
}
{}{}{.}

 }

__NOEDITSECTION__

\inputaufgabeklausurloesung
{2}

{

Wir betrachten die Standardparabel

\mavergleichskette
{\vergleichskette
{ Y
}
{ \subseteq }{ \R^2
}
{  }{
}
{    }{
}
{   }{
}
}
{}{}{}
mit der induzierten riemannschen Struktur und mit der Parametrisierung
\maabbeledisp {\varphi} {\R} { Y
} { t } { \left(  t   , \,  t^2                   \right)
} {,}
die wir als eine
\zusatzklammer {inverse} {} {}
Karte betrachten. Bestimme die riemannsche
\definitionsverweis {Fundamentalfunktion}{}{}

\mavergleichskette
{\vergleichskette
{ g(t)
}
{ = }{ g_{11} (t)
}
{  }{
}
{    }{
}
{   }{
}
}
{}{}{.}

}
{

Es ist

\mavergleichskettealign
{\vergleichskettealign
{ g(t)
}
{ =} { \left\langle  T_t \varphi(1)   , T_t \varphi(1)   \right\rangle_{ Y, \varphi(t) }
}
{ =} { \left\langle    \varphi' (t)   ,   \varphi'(t)   \right\rangle
}
{ =} { \left\langle  \begin{pmatrix}  1 \\ 2t   \end{pmatrix}     ,  \begin{pmatrix}  1 \\ 2t   \end{pmatrix}   \right\rangle
}
{ =} { 1+4t^2
}
}
{}
{}{.}

 }

__NOEDITSECTION__

\inputaufgabeklausurloesung
{6}

{

Berechne die zurückgezogene Differentialform
\mathl{\varphi^* \tau}{} zu

\mavergleichskettedisp
{\vergleichskette
{ \tau
}
{ =} { dx \wedge dy \wedge dz - w dx \wedge dy \wedge dw + \cos (xy)  dx \wedge dz \wedge dw-yw dy \wedge dz \wedge dw
}
{ } {
}
{ } {
}
{ } {
}
}
{}{}{}
unter der Abbildung
\maabbeledisp {\varphi} { \R^3 } { \R^4
} { (r,s,t) } {  \left(  r^2s,t, \sin  r ,e^{st}  \right)  = (x,y,z,w)
} {.}

}
{

Es ist

\mavergleichskettedisp
{\vergleichskette
{dx
}
{ =} { d(r^2s)
}
{ =} { 2rs dr  +r^2ds
}
{ } {
}
{ } {
}
}
{}{}{,}

\mavergleichskettedisp
{\vergleichskette
{dy
}
{ =} { dt
}
{ } {
}
{ } {
}
{ } {
}
}
{}{}{,}

\mavergleichskettedisp
{\vergleichskette
{dz
}
{ =} { d (\sin  r )
}
{ =} { \cos  r  dr
}
{ } {
}
{ } {
}
}
{}{}{}
und

\mavergleichskettedisp
{\vergleichskette
{dw
}
{ =} { d(e^{st})
}
{ =} { se^{st} dt+te^{st} ds
}
{ } {
}
{ } {
}
}
{}{}{.}
Damit ist

\mavergleichskettealign
{\vergleichskettealign
{ \varphi^* \tau
}
{ =} { d(r^2s) \wedge dt \wedge d( \sin  r )  -e^{st} d(r^2s) \wedge dt \wedge d(e^{st})
}
{ \, \, \, \,} {+ \cos (r^2st)  d (r^2s) \wedge d ( \sin  r ) \wedge d(e^{st}) - te^{st} dt \wedge d( \sin  r ) \wedge d(e^{st})
}
{ =} { r^2 \cos (r^2st) \cos  r  ds \wedge dt \wedge dr  - 2rs t e^{st} e^{st} dr \wedge dt \wedge ds
}
{ \, \, \, \,} { + r^2 se^{st} \cos  r  ds \wedge dr \wedge dt  -  t^2 e^{st}  e^{st} \cos  r dt \wedge dr \wedge ds
}
}
{
\vergleichskettefortsetzungalign
{ =} { (r^2 \cos  r +2 rst e^{2st} - r^2 s e ^{st} \cos (r^2st) \cos  r - t^2e^{ 2st} \cos  r     )  dr \wedge ds \wedge dt
}
{ } {}
{ } {}
{ } {}
}
{}{}

 }

__NOEDITSECTION__

\inputaufgabeklausurloesung
{7 (2+3+2)}

{

Wir betrachten die differenzierbaren Abbildungen
\maabbeledisp {\gamma} { [1,2] } {\R^2
} { t } {(t,t^{-1})
} {,}
und
\maabbeledisp {\varphi} {\R^2} {\R^3
} {(u,v)} {(u^2,uv,-u+v^2)
} {,}
und die Differentialform

\mavergleichskettedisp
{\vergleichskette
{ \omega
}
{ =} { xdx-zdy+dz
}
{ } {
}
{ } {
}
{ } {
}
}
{}{}{}
auf dem $\R^3$.

a) Berechne die zurückgezogene Differentialform
\mathl{\varphi^* \omega}{} auf dem $\R^2$.

b) Berechne das Wegintegral zur Differentialform
\mathl{\varphi^* \omega}{} zum Weg $\gamma$.

c) Berechne
\zusatzklammer {ohne Bezug auf b)} {} {}
das Wegintegral zur Differentialform
\mathl{\omega}{} zum Weg $\varphi \circ \gamma$.

}
{

a) Die zurückgezogene Differentialform ist

\mavergleichskettealign
{\vergleichskettealign
{\varphi^* ( xdx -zdy +dz )
}
{ =} { u^2 du^2    -(-u+v^2) duv +d (-u+v^2)
}
{ =} { 2u^3 du + (u-v^2) (vdu +udv) -du +2vdv
}
{ =} { (2u^3 +uv -v^3 -1) du    + (u^2 -uv^2 +2v )dv
}
{ } {
}
}
{}
{}{.}

b) Das Wegintegral ist

\mavergleichskettealign
{\vergleichskettealign
{ \int_1^2 ( 2t^3 +1 -t^{-3} -1 )dt  +(t^2 -t^{-1} +2t^{-1} ) d t^{-1}
}
{ =} { \int_1^2 ( 2t^3 -t^{-3} ) dt  +(t^2 + t^{-1} ) d t^{-1}
}
{ =} { \int_1^2 ( 2t^3 -t^{-3} ) dt  -(1 + t^{-3} ) dt
}
{ =} { \int_1^2 ( 2t^3 - 2t^{-3} -1 ) dt
}
{ =} { \left(  { \frac{   1 }{  2 } }  t^4  + t^{-2}  -t  \right) \vert_{  1 }^{  2 }
}
}
{
\vergleichskettefortsetzungalign
{ =} { 8 + { \frac{   1 }{  4 } }  -2 - { \frac{   1 }{  2 } } -1 +1
}
{ =} { 5 + { \frac{   3 }{  4 } }
}
{ } {}
{ } {}
}
{}{.}

c) Der verknüpfte Weg ist

\mavergleichskettedisp
{\vergleichskette
{ \psi (t)
}
{ =} { (t^2,1,-t +t^{-2} )
}
{ } {
}
{ } {
}
{ } {
}
}
{}{}{.}
Somit ist

\mavergleichskettealign
{\vergleichskettealign
{ \int_1^2 \psi^* (xdx -zdy +dz )
}
{ =} { \int_1^2 t^2dt^2 + d (-t + t^{-2} )
}
{ =} {  \int_1^2 (2 t^3-1 -2t^{-3} ) dt
}
{ =} { 5 + { \frac{   3 }{  4 } }
}
{ } {
}
}
{}
{}{.}

 }

__NOEDITSECTION__

\inputaufgabeklausurloesung
{4}

{

Es sei $M$ eine
\definitionsverweis {Mannigfaltigkeit mit Rand}{}{.}
Zeige, dass jede offene Teilmenge

\mavergleichskette
{\vergleichskette
{ N
}
{ \subseteq }{ M
}
{  }{
}
{    }{
}
{   }{
}
}
{}{}{}
ebenfalls eine Mannigfaltigkeit mit
\zusatzklammer {eventuell leerem} {} {}
Rand ist, und dass

\mavergleichskettedisp
{\vergleichskette
{ \partial N
}
{ =} { \partial M \cap N
}
{ } {
}
{ } {
}
{ } {
}
}
{}{}{}
gilt.

}
{

Es sei
\mathl{P \in N}{} und
\mathl{P \in U}{} ein Kartengebiet von $M$. Es sei
\maabbdisp {\alpha} { U } {V
} {}
eine Karte mit einer offenen Menge
\mathl{V \subseteq H}{,}
\mathl{H=\R_{\geq 0} \times \R^{n-1}}{.} Diese Karte induziert einen Homöomorphismus
\maabbdisp {} {N \cap U} { \alpha (N \cap U)
} {.}
Diese Karten nehmen wir für $N$. Die Diffeomorphieeigenschaft der Kartenwechsel zu $M$ überträgt sich direkt auf die Karten zu $N$. Es sei
\mathl{P \in N \cap \partial M}{.} Dann wird $P$ unter einer Karte auf einen Randpunkt von $H$ abgebildet, und da die Karten für $N$ Einschränkungen von solchen Karten sind, bleibt diese Eigenschaft erhalten. Wenn umgekehrt
\mathl{P \in \partial N}{} ist, so ist natürlich einerseits
\mathl{P \in N \subseteq M}{.} Andererseits gibt es zu $P$ eine Karte zu $N$, die durch Einschränkung von einer Karte zu $M$ herrührt, in der $P$ auf einen Randpunkt von $H$ abgebildet wird.

 }

__NOEDITSECTION__

\inputaufgabeklausurloesung
{8}

{

Beweise den Satz über die Charakterisierung von geodätischen Kurven bezüglich des Levi-Civita-Zusammenhangs.

}
{

Es sei $n$ die Dimension von $M$. Wir betrachten die Situation direkt auf einem offenen Kartenbild

\mavergleichskette
{\vergleichskette
{ U
}
{ \subseteq }{ \R^n
}
{  }{
}
{    }{
}
{   }{
}
}
{}{}{.}
Die vertikale Ableitung ist gemäß
[[Offene Menge/Triviales Vektorbündel/Linearer Zusammenhang/Spaltung/Christoffelsymbole/Bemerkung|Bemerkung]] [[Kurs:Differentialgeometrie/Offene Menge/Triviales Vektorbündel/Linearer Zusammenhang/Spaltung/Christoffelsymbole/Bemerkung/Bemerkungreferenznummer|*****]]
durch
\maabbeledisp {} {T(U \times \R^n) = U \times \R^n \times \R^n \times \R^n } { U \times \R^n \times  \R^n
} {(P,e_j, \partial_i, w)} {(P,e_j,\sum_{k = 1}^n \Gamma^k_{ij} (P) e_k +w) \cong V(U \times \R^n)
} {}
gegeben, wobei man die Abbildung nach $TU$ erhält, wenn man die mittlere Komponente weglässt. Die zweite Ableitung der Kurve ist zunächst die zweite Tangentialabbildung, es ist
\zusatzklammer {wobei wir die Multiplikation mit der eindimensionalen Richtung des Tangentialraumes der Kurve ignorieren} {} {}
\maabbeledisp {T(\gamma)} { I } { TU
} { t } { (\gamma(t), \gamma'(t))
} {,}
und entsprechend
\maabbeledisp {T(T(\gamma))} { I } { T(TU)
} { t } { ((\gamma(t), \gamma'(t)), (\gamma'(t), \gamma^{\prime \prime} (t)))
} {}
\zusatzklammer {es werden beide Komponenten der Tangentialabbildung abgeleitet} {} {.}
Mit

\mavergleichskette
{\vergleichskette
{ \gamma'(t)
}
{ = }{ \sum_{j = 1}^n \gamma_j'(t) \partial_j
}
{  }{
}
{    }{
}
{   }{
}
}
{}{}{}
wird dies unter der vertikalen Projektion auf

\mathdisp {\sum_{k = 1}^n  \left(  \sum_{i,j} \gamma_i'(t) \gamma_j'(t) \Gamma^k_{ij}(\gamma(t))  \right)  \partial_k  +\sum_{k = 1}^n \gamma_k^{\prime \prime}(t) \partial_k} {  }

abgebildet. Die Bedingung an eine Geodätische, dass die zweite Ableitung
\zusatzklammer {in $TTM$} {} {}
stets horizontal ist, ist äquivalent dazu, dass die berechnete vertikale Projektion gleich $0$ ist. Dies bedeutet, dass die einzelnen Komponenten gleich $0$ sind und dies bedeutet

\mavergleichskettedisp
{\vergleichskette
{ \sum_{i,j} \gamma_i'(t) \gamma_j'(t) \Gamma^k_{ij}(\gamma(t))  + \gamma_k^{\prime \prime}(t)
}
{ =} { 0
}
{ } {
}
{ } {
}
{ } {
}
}
{}{}{}
für

\mavergleichskette
{\vergleichskette
{ k
}
{ = }{ 1  , \ldots , n
}
{  }{
}
{    }{
}
{   }{
}
}
{}{}{.}

 }

 [[Kategorie:Latexseite]]
